% ===================================================================
%  CHART No. 5 — The House and how it is built.
%  Traduction 2026 d'apres texto/io, controlee sur texto/fr.
%  Consigne : voir texto/en/10-table-01.tex.
%
%  Le tableau 5 est un DIALOGUE, et les deux volumes composent le nom
%  du parleur en petites capitales : on garde \textsc{}, que le rendu
%  laisse tomber en gardant le texte.
%
%  Les renvois du plan sont des LETTRES — (a) le couloir, (b) le
%  salon... — et non des numeros. On les ecrit comme l'ido les ecrit,
%  y compris la ou l'imprimeur a compose « (1) » pour « (l) » : la
%  page de lecture a deja de quoi trancher, et l'anglais n'a pas a
%  corriger de son cote ce que les deux autres colonnes portent.
% ===================================================================

%%K t05-apar-1 apar
\VUsaut{6.0mm}
\VUtitre{15.0pt}{60}{CHART No. 5}
\VUsaut{1.4mm}
\VUtitre{11.4pt}{0}{The House and how it is built.}
\VUsaut{2.2mm}

%%K t05-tit-1 sub
\VUpk{10.2pt}{40}{A lucky meeting.}
\VUsaut{3.0mm}

%%K t05-01-1 p
\textsc{John}. --- Uncle, I would really like to know how a
\VUgras{house} is built.

%%K t05-01-2 p
\textsc{The Uncle} \textsuperscript{(80)}. --- That's easy enough, my boy. Come with me and
I'll show you the one Mr Bernard \textsuperscript{(79)} is having built, the one I'm going to
live in.

%%K t05-01-3 p
\textsc{John}. --- And who drew the \VUgras{plan} \textsuperscript{(76)} of this house?

%%K t05-01-4 p
\textsc{The Uncle}. --- The \VUgras{architect} \textsuperscript{(75)}; he drew it up very
clearly, it was approved, and the \VUgras{building contractor} \textsuperscript{(77)} started
work.

%%K t05-01-5 p
\textsc{John}. --- And who is the contractor?

%%K t05-01-6 p
\textsc{The Uncle}. --- Mr White, the father of your friend James. He runs the site along with
Mr Roux, the architect.

%%K t05-01-7 p
Well, look at that! Here comes Mr White at just the right moment with his
\VUgras{rule} \textsuperscript{(78)}, and Mr Roux with his plan.

%%K t05-01-8 p
Good morning, gentlemen. How are you?

%%K t05-01-9 p
\textsc{Mr Roux}. --- Very well, thanks. Morning, John --- hasn't he grown!

%%K t05-01-10 p
\textsc{The Uncle}. --- He has, hasn't he. He's quite the little man now. He is very serious;
everything he sees has to be explained to him. Today he would like to know how a house is
built.

%%K t05-tit-2 sub
\VUpk{10.2pt}{40}{The workmen.}
\VUsaut{3.0mm}

%%K t05-01-11 p
\textsc{Mr White}. --- Come along with me, then, and I'll explain it to you right now. I do
like children who want to learn.

%%K t05-01-12 p
You see that workman with a \VUgras{shovel} \textsuperscript{(82)} in his hand: he is a
\VUgras{navvy} \textsuperscript{(81)}. He is digging a \VUgras{trench} \textsuperscript{(46)}
to lay the water pipe. He also dug out the ground where the house stands, so that the
bricklayer could put up the four
\VUgras{walls}. The part of those walls that is in the ground is called
the \VUgras{foundations} \textsuperscript{(2)}. The space enclosed by the foundations forms
the \VUgras{cellar} \textsuperscript{(3)}. This cellar has a small window called an
\VUgras{air vent} \textsuperscript{(5)}, a \VUgras{door} \textsuperscript{(4)} and a stairway.

%%K t05-01-13 p
The \VUgras{bricklayer} \textsuperscript{(108)} went on raising the walls, leaving openings
for the \VUgras{doors} \textsuperscript{(8)} and the \VUgras{windows} \textsuperscript{(17)}.

%%K t05-01-14 p
\textsc{John}. --- But I can't see him, this bricklayer!

%%K t05-01-15 p
\textsc{Mr White}. --- He has finished working on the house now. He is over by the
\VUgras{stable} \textsuperscript{(54)}, up on his \VUgras{scaffolding}
\textsuperscript{(110)}, building the wall of the \VUgras{wash house} \textsuperscript{(56)}.

%%K t05-01-16 p
He has a trowel, a trough and a \VUgras{plumb line} \textsuperscript{(109)}. But you won't
remember those names.

%%K t05-01-17 p
\textsc{John}. --- Oh yes I will, because I'm going to ask my uncle for a little trowel and a
trough, and I'll make a plumb line myself with a piece of string.

%%K t05-tit-3 sub
\VUsaut{3.0mm}
\VUpk{10.2pt}{40}{The materials.}
\VUsaut{3.0mm}

%%K t05-01-18 p
\textsc{The Uncle}. --- Ah, very good. But tell me, John, what do you need to build a house?

%%K t05-01-19 p
\textsc{John}. --- You need \VUgras{dressed stone} \textsuperscript{(59)} and
\VUgras{mortar} \textsuperscript{(65)}.

%%K t05-01-20 p
\textsc{The Uncle}. --- Exactly. Look at that man in the big hat, on his
\VUgras{cart} \textsuperscript{(136)}. He is a \VUgras{carter} \textsuperscript{(135)}. Every
day he brings \VUgras{blocks of stone} \textsuperscript{(60)} here. He unloads his cart in the
yard, and the \VUgras{stonecutters} \textsuperscript{(83)} dress the blocks and saw them with
their \VUgras{stone saw} \textsuperscript{(85)}. A
\VUgras{labourer} \textsuperscript{(146)} carries them to the bricklayer, who sets
them in place.

%%K t05-01-21 p
Meanwhile the \VUgras{mortar mixer} \textsuperscript{(87)} is preparing the mortar. He needs
\VUgras{sand} \textsuperscript{(62)}, \VUgras{lime} \textsuperscript{(63)} and
\VUgras{water} \textsuperscript{(64)}. The lime is beside him in a \VUgras{sack}
\textsuperscript{(71)}; a \VUgras{bricklayer's mate} \textsuperscript{(111)} brings him the
sand in a \VUgras{wheelbarrow} \textsuperscript{(112)}, and the \VUgras{apprentice}
\textsuperscript{(113)} is at the pump \textsuperscript{(58)}. He is filling a
\VUgras{bucket} \textsuperscript{(114)}. The mortar will soon be ready. The
\VUgras{hod carrier} \textsuperscript{(107)} takes it and carries it up the ladder
to the scaffolding, where a bricklayer is finishing that side of the house.

%%K t05-01-22 p
And now, what do you call the workman who makes the \VUgras{roof} \textsuperscript{(31)}?

%%K t05-01-23 p
\textsc{John}. --- He is the \VUgras{roofer} \textsuperscript{(118)}.

%%K t05-tit-4 sub
\VUsaut{3.0mm}
\VUpk{10.2pt}{40}{The roof of the house.}
\VUsaut{3.0mm}

%%K t05-01-24 p
\textsc{Mr White}. --- Yes, but first comes the \VUgras{carpenter} \textsuperscript{(115)}.

%%K t05-01-25 p
The carpenter makes the \VUgras{roof frame} \textsuperscript{(35)}. He joins the
\VUgras{beams} \textsuperscript{(37)} with \VUgras{pegs} \textsuperscript{(117)}. Those
workmen up there, in their wide velvet breeches, are carpenters. One is holding a small beam
and the other is cutting a peg with his \VUgras{axe} \textsuperscript{(116)}. The one near the
\VUgras{chimney} \textsuperscript{(41)} is the roofer. He nails laths on to the
\VUgras{joists} (*), and \VUgras{slates} \textsuperscript{(66)} on to the laths. Sometimes he
lays tiles.

%%K t05-noto-1 noto
\VUnotes{143.2mm}{%
(*) \VUgras{Balk-o} = German \textit{Balken}, English \textit{joist}, French \textit{solive},
Italian \textit{travicello}, Russian \textit{balka}, Spanish and Portuguese \textit{viga}. A
technical root not yet adopted, but proposed here to render the sense of the French word
\textit{solive}, which is neither a \textit{trabo} (French \textit{poutre}) nor a small beam.
It is a squared timber that carries a floor and a ceiling.%
}

%%K t05-01-26 p
The roof of the stable is \VUgras{tiled} \textsuperscript{(55)}, that of the house is
\VUgras{slated} \textsuperscript{(32)}, and that of the veranda is of
\VUgras{zinc} \textsuperscript{(24)}.

%%K t05-01-27 p
\textsc{John}. --- Does the roofer make the zinc roofs as well?

%%K t05-01-28 p
\textsc{Mr White}. --- No, that's the \VUgras{zinc worker} \textsuperscript{(121)} or the
\VUgras{plumber}, the one setting a \VUgras{dormer window} \textsuperscript{(38)} into the
roof of the house. He also lays the \VUgras{roof gutters} \textsuperscript{(43)} and the
\VUgras{downpipes} \textsuperscript{(42)}. He solders the zinc and the tinplate. He is the one
who fixed the \VUgras{lightning conductor} \textsuperscript{(40)} and the \VUgras{weathervane}
\textsuperscript{(39)} on the
\VUgras{gable} \textsuperscript{(36)}.

%%K t05-01-29 p
\textsc{John}. --- Is the house finished, then?

%%K t05-tit-5 sub
\VUsaut{3.0mm}
\VUpk{10.2pt}{40}{The tools.}
\VUsaut{3.0mm}

%%K t05-01-30 p
\textsc{Mr White}. --- Not quite. The \VUgras{plasterer} \textsuperscript{(99)} builds with
\VUgras{bricks} \textsuperscript{(61)} the partitions that divide it into the different rooms.
He lays \VUgras{plaster} \textsuperscript{(70)} on the
\VUgras{ceilings} \textsuperscript{(26)} and on the walls. Just now he is working on the ceiling
of the \VUgras{veranda} \textsuperscript{(33)}. Like the bricklayer, he has a
\VUgras{trowel} \textsuperscript{(100)} and a \VUgras{trough} \textsuperscript{(101)}.

%%K t05-01-31 p
Then the joiner makes the doors, the windows, the \VUgras{wooden floors}
\textsuperscript{(16)} and the \VUgras{shutters} \textsuperscript{(19)}.

%%K t05-01-32 p
Just now he is working in the yard at his \VUgras{workbench} \textsuperscript{(90)}. He has a
\VUgras{saw} \textsuperscript{(94)} to cut the \VUgras{wood} \textsuperscript{(69)}, a
\VUgras{plane} \textsuperscript{(91)}, a \VUgras{holdfast} \textsuperscript{(92)}, a
\VUgras{chisel} \textsuperscript{(94 \textit{bis})}, a
\VUgras{pair of compasses} \textsuperscript{(95)} and a wooden \VUgras{mallet}
\textsuperscript{(95 \textit{bis})}.

%%K t05-01-33 p
\textsc{John}. --- Oh, but being a joiner is even more fun than being a bricklayer. Uncle,
will you buy me a workbench, a saw and a plane? I'm going to make a kennel for Medor.

%%K t05-01-34 p
\textsc{The Uncle}. --- Yes, my boy.

%%K t05-01-35 p
\textsc{John}. --- Oh, thank you! And who's that man standing near the front door?

%%K t05-tit-6 sub
\VUsaut{3.0mm}
\VUpk{10.2pt}{40}{The painting.}
\VUsaut{3.0mm}

%%K t05-01-36 p
\textsc{Mr White}. --- He is a \VUgras{painter} \textsuperscript{(102)}. He is standing on his
ladder with a pot of \VUgras{paint} \textsuperscript{(103)} and his
\VUgras{brush} \textsuperscript{(104)}. He is coating the wood of the front door so
that it will last longer.

%%K t05-01-37 p
He is also the one who hangs the wallpaper in the rooms.

%%K t05-01-38 p
The \VUgras{upholsterer} \textsuperscript{(107)} on a \VUgras{ladder} \textsuperscript{(106)},
out on the \VUgras{balcony} \textsuperscript{(28)}, is putting up a \VUgras{blind}
\textsuperscript{(29)}.

%%K t05-01-39 p
The workman standing by the big window is the \VUgras{glazier} \textsuperscript{(96)}. He
fixes the \VUgras{panes} \textsuperscript{(18)} with
\VUgras{putty} \textsuperscript{(73)}. That other workman, kneeling by the cellar
door, is a \VUgras{locksmith} \textsuperscript{(97)}. He is fitting a \VUgras{lock}
\textsuperscript{(6)}. His \VUgras{file} \textsuperscript{(98)} is beside him.

%%K t05-01-40 p
\textsc{John}. --- Is the man hitting a piece of iron with his
\VUgras{hammer} \textsuperscript{(84)} a locksmith too?

%%K t05-01-41 p
\textsc{Mr White}. --- Not quite --- he's a \VUgras{blacksmith} \textsuperscript{(125)}. He is
forging \VUgras{ironwork} \textsuperscript{(67)} for the
\VUgras{electrician} \textsuperscript{(124)}, who is on the roof of the
\VUgras{coach house} \textsuperscript{(51)}. He has a \VUgras{forge} \textsuperscript{(126)},
an \VUgras{anvil} \textsuperscript{(127)} and a pair of \VUgras{tongs}
\textsuperscript{(128)}.

%%K t05-01-42 p
The electrician is fixing the \VUgras{copper wire} \textsuperscript{(44)} of the telephone.

%%K t05-tit-7 sub
\VUpk{10.2pt}{40}{The gardening.}
\VUsaut{3.0mm}

%%K t05-01-43 p
\textsc{The Uncle}. --- At the far end of the \VUgras{yard} \textsuperscript{(45)}, near the
\VUgras{railings} \textsuperscript{(47)} and the \VUgras{carriage gate}
\textsuperscript{(48)}, you can see the \VUgras{gardener} \textsuperscript{(129)}. He is
getting the little \VUgras{garden} \textsuperscript{(49)} ready. He has dug the soil over with
his \VUgras{spade} \textsuperscript{(131)}, he is sowing seed and raking with the
\VUgras{rake} \textsuperscript{(130)}. Then, with his
\VUgras{watering can} \textsuperscript{(132)}, he will water the \VUgras{flower beds}
\textsuperscript{(150)} and the plants will grow.

%%K t05-01-44 p
The \VUgras{groom} \textsuperscript{(133)}, who has just tended the horse and fed it, is
cleaning the \VUgras{carriage} \textsuperscript{(52)} with a sponge, a
\VUgras{hand pump} \textsuperscript{(134)} and a bucket of water. Then he will put
it away in the coach house and shut the door with the heavy \VUgras{bolt}
\textsuperscript{(53)}.

%%K t05-01-45 p
There, I hope you are satisfied now --- you have had plenty of explanations.

%%K t05-01-46 p
\textsc{John}. --- Yes, uncle, but I'd like to see the inside of the house.

%%K t05-01-47 p
\textsc{The Uncle}. --- That can wait for tomorrow. Mr Roux will explain the plan to you and
tell you the name of each room.

%%K t05-tit-8 sub
\VUsaut{3.0mm}
\VUpk{10.2pt}{40}{The layout of the house.}
\VUsaut{2.4mm}

%%K t05-apar-2 apar
{\centering\textit{(See the plan.)}\par}
\VUsaut{2.4mm}

%%K t05-01-48 p
\textsc{The Architect}. --- Come along, John, I'm going to take you round the different parts
of the house.

%%K t05-01-49 p
Let's go in on the \VUgras{ground floor} \textsuperscript{(7)}. Here is the
\VUgras{doorbell} button \textsuperscript{(11)}. We climb the three \VUgras{steps}
\textsuperscript{(14)} of the \VUgras{front steps} \textsuperscript{(9)}, we cross the
\VUgras{threshold} \textsuperscript{(10)} of the \VUgras{front door} \textsuperscript{(8)},
and here we are at the foot of the \VUgras{staircase} \textsuperscript{(13)}.

%%K t05-01-50 p
\textsc{John}. --- This is the corridor.

%%K t05-01-51 p
\textsc{The Architect}. --- No, this is the \VUgras{hall} \textsuperscript{(12)}. The
\VUgras{corridor} \textsuperscript{(a)} is at the far end. On the right, here are
the \VUgras{drawing room} \textsuperscript{(b)} and the \VUgras{dining room}
\textsuperscript{(c)}; opposite the dining room are the \VUgras{kitchen} \textsuperscript{(d)}
and the \VUgras{pantry} \textsuperscript{(e)}; then, at the front, the \VUgras{study}
\textsuperscript{(f)}. The \VUgras{veranda} \textsuperscript{(23)}, with its
\VUgras{glazed door} \textsuperscript{(25)}, is next to the drawing room.

%%K t05-01-52 p
Now we'll go up the stairs. Hold on to the \VUgras{handrail} \textsuperscript{(15)} and count
the steps. There are fourteen. Here is the
\VUgras{landing} \textsuperscript{(m)}: we are on the first \VUgras{floor}
\textsuperscript{(27)}.

%%K t05-tit-9 sub
\VUsaut{3.0mm}
\VUpk{10.2pt}{40}{The first floor.}
\VUsaut{3.0mm}

%%K t05-01-53 p
Opposite the stairs are the \VUgras{lavatory} \textsuperscript{(20)} and the
\VUgras{dressing room} \textsuperscript{(j)}. On the right, a fine
\VUgras{bedroom} \textsuperscript{(g)} with two windows on the \VUgras{front}
\textsuperscript{(1)} of the house. On the left are the \VUgras{bathroom}
\textsuperscript{(1)} and the \VUgras{guest room} \textsuperscript{(i)} with a large
\VUgras{alcove} \textsuperscript{(h)}. Next to the bedroom is the
\VUgras{children's room} \textsuperscript{(k)}. It opens on to a wide
\VUgras{terrace} \textsuperscript{(21)}. Under this terrace is another lavatory
\textsuperscript{(20)}, and against the wall, here is the \VUgras{bell} \textsuperscript{(22)}
that is rung for dinner.

%%K t05-01-54 p
\textsc{John}. --- Let's go up to the \VUgras{second floor}.

%%K t05-01-55 p
\textsc{The Architect}. --- There is no second floor. Under the roof there are only
\VUgras{attic rooms} \textsuperscript{(33)} and the loft \textsuperscript{(34)}. That is the
whole tour --- we can go back down.

%%K t05-01-56 p
\textsc{John}. --- Thank you, sir, that was really interesting. I'm going to tell my father
everything I have seen.
\VUsaut{4.0mm}
\VUfilet{22mm}
